\documentclass[10pt]{article}

\usepackage[utf8]{inputenc}
\usepackage[T1]{fontenc}
\usepackage{amsmath,amssymb,mathtools}
\usepackage{booktabs}
\usepackage{enumitem}
\usepackage[margin=0.78in]{geometry}
\usepackage{xcolor}
\usepackage[most]{tcolorbox}
\usepackage{hyperref}

\definecolor{BootcampBlue}{RGB}{31,78,121}
\definecolor{BootcampGold}{RGB}{215,155,35}
\definecolor{BootcampGreen}{RGB}{34,105,72}
\hypersetup{colorlinks=true,linkcolor=BootcampBlue,urlcolor=BootcampBlue}
\setlist{itemsep=1pt,topsep=2pt,parsep=0pt,partopsep=0pt}
\setlength{\parindent}{0pt}
\setlength{\parskip}{3pt}
\setlength{\abovedisplayskip}{5pt plus 2pt minus 2pt}
\setlength{\belowdisplayskip}{5pt plus 2pt minus 2pt}

\newcommand{\R}{\mathbb{R}}
\newcommand{\E}{\mathbb{E}}
\newcommand{\Var}{\operatorname{Var}}
\newcommand{\Cov}{\operatorname{Cov}}
\newcommand{\trans}{^{\mathsf T}}
\newcommand{\norm}[1]{\left\lVert#1\right\rVert}
\newcommand{\diag}{\operatorname{diag}}

\newtcolorbox{definitionbox}[1]{breakable,colback=BootcampBlue!4,
  colframe=BootcampBlue,boxrule=0.8pt,arc=1mm,
  title=\textbf{#1},fonttitle=\bfseries}
\newtcolorbox{checkpointbox}[1]{breakable,colback=BootcampGreen!5,
  colframe=BootcampGreen,boxrule=0.8pt,arc=1mm,
  title=\textbf{#1},fonttitle=\bfseries}
\newtcolorbox{deliverablebox}[1]{breakable,colback=BootcampGold!8,
  colframe=BootcampGold!80!black,boxrule=0.8pt,arc=1mm,
  title=\textbf{#1},fonttitle=\bfseries}

\title{Summer Bootcamp -- Session 3\\Portfolio Optimization Project}
\author{Nathan Sauldubois}
\date{August 2026}

\begin{document}
\maketitle

\section*{Optional project}

This optional extension can be completed independently by students who want to
explore portfolio optimization further. The aim is to build a complete and
reproducible research pipeline: simulate train and test returns, study the
covariance structure with PCA, estimate several portfolios, and evaluate
rebalanced strategies without look-ahead bias. The central question is:

\begin{quote}
Does a better-conditioned covariance estimate produce more stable weights and
better realized risk after transaction costs?
\end{quote}

The stages and self-checks provide a suggested path that students can follow at
their own pace. They may stop after any stage or select the parts that interest
them. A partial pipeline with correct validation is more useful than a more
sophisticated strategy containing look-ahead bias.

\begin{deliverablebox}{Files supplied and suggested outputs}
The folder contains this statement and
\texttt{portfolio\_project\_skeleton.py}. A complete implementation can create
\texttt{returns\_train.csv}, \texttt{returns\_test.csv}, and
\texttt{true\_parameters.npz}. Store tables and figures in clearly named
output folders. Do not edit the generated CSV files manually.
\end{deliverablebox}

\section{Definitions and conventions}

\begin{definitionbox}{Returns, moments, and annualization}
For $n$ assets, $R_t=(R_{1,t},\ldots,R_{n,t})\trans$ is the random vector of
daily simple returns. A sample contains rows indexed by date and one column per
asset. The population daily moments are
\[
\mu_d=\E[R_t],\qquad \Sigma_d=\Cov(R_t).
\]
Under the usual independent-return annualization convention,
\[
\mu=252\mu_d,\qquad \Sigma=252\Sigma_d.
\]
This convention is an approximation for arithmetic returns; state it in the
report rather than treating it as an identity for compounded wealth.
\end{definitionbox}

\begin{definitionbox}{Portfolio return, risk, and constraints}
A portfolio is a vector $w\in\R^n$. Full investment means
$\mathbf1\trans w=1$; long-only means $w_i\geq0$. Its expected return and
variance are
\[
R(w)=\mu\trans w,\qquad \sigma^2(w)=w\trans\Sigma w.
\]
Equal weight uses $w_i=1/n$. The Herfindahl concentration index is
$H(w)=\sum_iw_i^2$; under long-only full investment it ranges from $1/n$ to
$1$.
\end{definitionbox}

\begin{definitionbox}{Portfolio optimization problems}
The global minimum-variance portfolio (GMV) solves
\[
\min_w\ \frac12w\trans\Sigma w
\quad\text{s.t.}\quad \mathbf1\trans w=1
\]
with any additional bounds stated in the question. A mean--variance portfolio
solves
\[
\min_w\ \frac12w\trans\Sigma w-\gamma\mu\trans w,
\quad \gamma\geq0.
\]
For annual risk-free rate $r_f$, the Sharpe ratio is
\[
\operatorname{SR}(w)=
\frac{\mu\trans w-r_f}{\sqrt{w\trans\Sigma w}}.
\]
A maximum-Sharpe portfolio maximizes this ratio over its feasible set. If the
denominator is numerically zero, the candidate is invalid.
\end{definitionbox}

\begin{definitionbox}{Covariance estimators used in the project}
Let $S=U\diag(\lambda_1,\ldots,\lambda_n)U\trans$, with decreasing
eigenvalues. The rank-$k$ PCA reconstruction is
\[
S_{\mathrm{PCA}}^{(k)}=U_k\diag(\lambda_1,\ldots,\lambda_k)U_k\trans.
\]
Use the diagonal-corrected version
\[
\widetilde S_{\mathrm{PCA}}^{(k)}=S_{\mathrm{PCA}}^{(k)}+
\diag\!\left(\diag(S-S_{\mathrm{PCA}}^{(k)})\right)
\]
for optimization. Ridge regularization uses
\[
S_\delta=S+\delta I,
\qquad \delta=\tau\,\operatorname{tr}(S)/n.
\]
The scale-free parameter $\tau$ makes experiments easier to compare.
\end{definitionbox}

\section{Suggested Python architecture}

Use the supplied skeleton and implement the functions below in order. Keep the
function names and return types so that later stages can reuse earlier work.
Each function should validate its own output before returning it.

\begingroup
\raggedright
\begin{enumerate}[label=\textbf{A.\arabic*.}]
\item \textbf{Data generation.} Implement
\path{factor_covariance},
\path{generate_population_parameters},
\path{theoretical_moments},
\path{simulate_returns}, and
\path{generate_and_save_data}.
The final function should return two data frames and create the two CSV files and
parameter archive. Add assertions for shapes, dates, column names, finite
values, and covariance eigenvalues.

\item \textbf{Loading and descriptive checks.} Implement
\path{load_project_data} and
\path{annualized_sample_moments}.
The loader should parse dates, reject overlapping train/test indices, preserve
asset order, and compare stored dimensions with the configuration.

\item \textbf{PCA.} Implement
\path{fit_pca_on_train},
\path{pca_covariance}, \path{pca_scores},
\path{principal_angles}, and
\path{plot_pca_diagnostics}.
Return eigenvalues in decreasing order and use the training mean and loadings
when transforming test observations.

\item \textbf{Portfolio construction.} Implement
\path{ridge_covariance}, \path{unconstrained_gmv},
\path{solve_long_only_gmv},
\path{solve_max_sharpe},
\path{solve_mean_variance}, and
\path{trace_efficient_frontier}.
Every numerical solver should return both weights and its raw solver result.
Never silently accept a failed solver.

\item \textbf{Optimality and feasibility checks.} Implement
\path{portfolio_diagnostics} and
\path{directional_optimality_check}.
At minimum report budget error, bound violations, predicted variance,
concentration, solver status, and whether small feasible perturbations reduce
the objective.

\item \textbf{Model selection.} Implement
\path{select_covariance_hyperparameters}.
Its output should contain one row per candidate, the exact estimation and
validation dates, realized validation risk, concentration, and the declared
selection criterion. The function should not accept test data as an argument.

\item \textbf{Backtesting.} Implement
\path{drift_weights}, \path{frozen_weight_backtest},
\path{backtest_rebalanced_strategy}, and
\path{performance_metrics}.
Return aligned series or data frames for target weights, pre-trade weights,
gross returns, net returns, turnover, costs, wealth, and drawdown.

\item \textbf{Reporting and robustness.} Implement
\path{plot_backtest_diagnostics} and \path{run_one_seed}.
The top-level \path{main} function can call the stages in order and
regenerate all files from an empty output folder without manual intervention.
\end{enumerate}
\endgroup

\begin{checkpointbox}{Recommended coding structure}
Avoid putting the full project in one notebook cell or one function. Each
optimizer, estimator, validation rule, metric, and plot should be callable and
testable independently. Use informative exceptions for invalid inputs and
store solver diagnostics alongside portfolio weights.
\end{checkpointbox}

\section{Detailed function contracts}

The following contracts describe a complete implementation. Shapes use
$T$ for the number of observations, $n$ for assets, $k$ for retained
components, and $m$ for equality constraints. Unless specified otherwise,
numeric arrays must be finite and use floating-point values.

\begin{sloppypar}
\subsection*{Utilities}

\paragraph{\texttt{make\_output\_directories(cfg)}}
\textbf{Input:} one \texttt{Config}. \textbf{Work:} create the data, figure, and
table directories with \texttt{parents=True} and \texttt{exist\_ok=True};
never delete existing results. \textbf{Return:} \texttt{None}.
\textbf{Check:} every expected path exists and is a directory.

\paragraph{\texttt{asset\_names(n\_assets)}}
\textbf{Input:} a positive integer. \textbf{Work:} generate zero-padded names
from \texttt{Asset\_01} through \texttt{Asset\_NN}.
\textbf{Return:} an ordered list of unique strings of length $n$.
\textbf{Reject:} nonpositive or noninteger input.

\paragraph{\texttt{check\_covariance(sigma, name, tol)}}
\textbf{Input:} a candidate matrix, a diagnostic name, and tolerance.
\textbf{Work:} verify square shape, finite entries, symmetry, and
$\lambda_{\min}\geq-\texttt{tol}$. \textbf{Return:} \texttt{None} on success.
\textbf{Failure:} raise an informative \texttt{ValueError} containing
\texttt{name}, the failed property, and the observed value.

\subsection*{Stage 1: simulation and files}

\paragraph{\texttt{factor\_covariance()}}
\textbf{Input:} none. \textbf{Work:} create $\mu_f$, factor volatilities, and
$R_f$, then compute
$\Sigma_f=\diag(s_f)R_f\diag(s_f)$ in annual units.
\textbf{Return:} three arrays with shapes $(3,)$, $(3,)$, and $(3,3)$.
\textbf{Checks:} correlation diagonal equals one, covariance diagonal equals
$s_f^2$, and covariance is positive semidefinite.

\paragraph{\texttt{generate\_population\_parameters(rng, cfg)}}
\textbf{Inputs:} a NumPy generator and configuration.
\textbf{Work:} draw $\alpha$, $B$, and idiosyncratic volatilities with the
stated distributions; build $D$; attach factor moments.
\textbf{Return:} a dictionary containing
\texttt{alpha}, \texttt{B}, \texttt{D}, \texttt{idio\_vol},
\texttt{mu\_f}, \texttt{factor\_vol}, and \texttt{sigma\_f}.
\textbf{Checks:} exact shapes, positive volatilities, diagonal $D$, and finite
values.

\paragraph{\texttt{theoretical\_moments(params)}}
\textbf{Input:} the population dictionary.
\textbf{Work:} compute annual
$\mu_{\mathrm{true}}=\alpha+B\mu_f$ and
$\Sigma_{\mathrm{true}}=B\Sigma_fB\trans+D$.
\textbf{Return:} arrays of shapes $(n,)$ and $(n,n)$.
\textbf{Checks:} expected keys and dimensions, symmetry, and positive
semidefiniteness.

\paragraph{\texttt{simulate\_returns(rng, n\_obs, params, cfg)}}
\textbf{Inputs:} generator, positive observation count, population parameters,
and configuration. \textbf{Work:} draw factor and idiosyncratic shocks in daily
units and combine them as
$\alpha/252+f_tB\trans+\varepsilon_t$.
\textbf{Return:} a $(T,n)$ array of daily simple returns.
\textbf{Checks:} exact shape, finite values, and no mutation of
\texttt{params}.

\paragraph{\texttt{returns\_frame(values, start\_date, cfg)}}
\textbf{Inputs:} a $(T,n)$ array and valid start date.
\textbf{Work:} attach an increasing business-day index and standard asset
names. \textbf{Return:} a data frame with named date index.
\textbf{Checks:} $n$ columns, unique dates, no missing values, and column order
equal to \texttt{asset\_names(n)}.

\paragraph{\texttt{generate\_and\_save\_data(cfg)}}
\textbf{Input:} configuration. \textbf{Work:} spawn independent RNG streams,
draw parameters once, simulate independent train and test shocks, create
non-overlapping dates, compute true moments, and save both CSV files plus the
NPZ archive. \textbf{Return:} train and test data frames.
\textbf{Checks:} reload the saved files, compare shapes and column order, and
verify that a second run with the same seed reproduces values exactly.

\paragraph{\texttt{load\_project\_data(cfg)}}
\textbf{Input:} configuration containing file paths.
\textbf{Work:} read CSV dates, load the NPZ archive, and reconstruct the
parameter dictionary. \textbf{Return:} train data frame, test data frame, and
parameter dictionary. \textbf{Reject:} missing files, wrong dimensions,
overlapping dates, reordered assets, nonnumeric returns, or missing parameter
keys.

\paragraph{\texttt{annualized\_sample\_moments(returns, cfg)}}
\textbf{Input:} one return data frame in daily units.
\textbf{Work:} compute $252\bar r$ and $252$ times the sample covariance with
the usual $T-1$ denominator. \textbf{Return:} $(n,)$ mean and $(n,n)$
covariance. \textbf{Checks:} at least two observations, all-numeric columns,
finite moments, and covariance validity.

\subsection*{Stage 2: PCA}

\paragraph{\texttt{fit\_pca\_on\_train(train)}}
\textbf{Input:} training returns only. \textbf{Work:} store the training mean,
center the data, compute daily sample covariance, apply
\texttt{numpy.linalg.eigh}, and sort eigenpairs decreasingly.
\textbf{Return:} a dictionary with \texttt{mean}, \texttt{eigenvalues},
\texttt{eigenvectors}, \texttt{explained\_ratio}, and
\texttt{empirical\_covariance}. \textbf{Checks:} centered means near zero,
orthonormal eigenvectors, nonincreasing eigenvalues, explained ratios summing
to one, and full reconstruction error below tolerance.

\paragraph{\texttt{pca\_covariance(eigenvalues, eigenvectors, n\_components,
empirical\_covariance, diagonal\_correction)}}
\textbf{Inputs:} fitted eigenpairs, integer $k$, empirical covariance, and
correction flag. \textbf{Work:} reconstruct the rank-$k$ covariance and,
when requested, add the diagonal residual.
\textbf{Return:} one symmetric $(n,n)$ covariance in the same units as the
input. \textbf{Checks:} $1\leq k\leq n$, compatible shapes, diagonal
preservation when corrected, and positive semidefiniteness.

\paragraph{\texttt{pca\_scores(returns, training\_mean,
training\_eigenvectors, n\_components)}}
\textbf{Inputs:} train or test returns and PCA quantities fitted on train.
\textbf{Work:} subtract the training mean and multiply by $U_k$; never refit or
recenter with the test mean. \textbf{Return:} a $(T,k)$ data frame named
\texttt{PC\_01}, \ldots, with the original dates.
\textbf{Checks:} matching asset dimension and valid $k$.

\paragraph{\texttt{principal\_angles(empirical\_vectors,
theoretical\_vectors, n\_components)}}
\textbf{Inputs:} two orthonormal bases and $k$.
\textbf{Work:} compute singular values of
$U_{\mathrm{true},k}\trans U_{\mathrm{emp},k}$, clip to $[-1,1]$, and take
arccos. \textbf{Return:} $k$ angles in radians, preferably sorted increasingly.
\textbf{Checks:} orthonormality and angles in $[0,\pi/2]$ up to tolerance.

\paragraph{\texttt{plot\_pca\_diagnostics(pca, output\_path)}}
\textbf{Inputs:} PCA dictionary and PNG path.
\textbf{Work:} make aligned scree and cumulative-explained-variance panels,
mark $k=3$, label units and axes, and call \texttt{tight\_layout}.
\textbf{Return:} \texttt{None}. \textbf{Checks:} create the parent directory,
save a nonempty file, and close the figure.

\subsection*{Stage 3: covariance estimators and portfolios}

\paragraph{\texttt{ridge\_covariance(sigma, tau)}}
\textbf{Inputs:} covariance and nonnegative $\tau$.
\textbf{Work:} compute
$\delta=\tau\operatorname{tr}(\Sigma)/n$ and return $\Sigma+\delta I$.
\textbf{Checks:} unchanged shape, symmetry, nondecreasing minimum eigenvalue,
and an informative error for $\tau<0$.

\paragraph{\texttt{equal\_weight(n\_assets)}}
\textbf{Input:} positive integer $n$.
\textbf{Return:} $(n,)$ vector with entries $1/n$.
\textbf{Checks:} finite nonnegative weights and budget error below tolerance.

\paragraph{\texttt{unconstrained\_gmv(sigma)}}
\textbf{Input:} positive-definite covariance.
\textbf{Work:} solve $\Sigma x=\mathbf1$ with
\texttt{numpy.linalg.solve} and normalize $w=x/(\mathbf1\trans x)$; do not form
an inverse. \textbf{Return:} $(n,)$ weights.
\textbf{Checks:} nonzero normalization, budget feasibility, and stationarity
$\Sigma w+\lambda\mathbf1\approx0$ for the implied multiplier.

\paragraph{\texttt{solve\_long\_only\_gmv(sigma, weight\_cap,
initial\_weights)}}
\textbf{Inputs:} covariance, cap, and optional feasible start.
\textbf{Work:} use SLSQP to minimize $\frac12w\trans\Sigma w$ with
$\mathbf1\trans w=1$ and $0\leq w_i\leq u$.
\textbf{Return:} weights and raw \texttt{OptimizeResult}.
\textbf{Reject:} $n\,u<1$, invalid covariance, infeasible start, or failed solver.
Here $n$ is the number of assets and $u$ the cap.

\paragraph{\texttt{solve\_max\_sharpe(mu, sigma, risk\_free\_rate,
weight\_cap, initial\_weights)}}
\textbf{Inputs:} annual moments, annual risk-free rate, bounds, and start.
\textbf{Work:} minimize negative Sharpe under the budget and bounds; evaluate
several feasible initializations outside the function.
\textbf{Return:} best successful weights and solver result.
\textbf{Checks:} positive predicted variance, finite objective, feasibility,
and no silent selection of a failed run.

\paragraph{\texttt{solve\_mean\_variance(mu, sigma, gamma, weight\_cap)}}
\textbf{Inputs:} annual moments, $\gamma\geq0$, and cap.
\textbf{Work:} minimize
$\frac12w\trans\Sigma w-\gamma\mu\trans w$ with analytical gradient, budget,
and bounds. \textbf{Return:} weights and solver result.
\textbf{Checks:} units consistent, feasible cap, solver success, and
constraints within tolerance.

\paragraph{\texttt{trace\_efficient\_frontier(mu, sigma, target\_returns,
weight\_cap)}}
\textbf{Inputs:} annual moments, ordered target grid, and cap.
\textbf{Work:} solve bounded minimum variance for every feasible target,
warm-starting from the preceding solution.
\textbf{Return:} one data-frame row per target containing requested and achieved
return, volatility, concentration, weights or a weight key, solver status,
message, and maximum violation. \textbf{Rule:} keep failed rows; do not hide
them from the report.

\paragraph{\texttt{portfolio\_diagnostics(weights, mu, sigma, lower, upper)}}
\textbf{Inputs:} one weight vector, annual moments, and bounds.
\textbf{Return:} a dictionary containing budget error, lower and upper
violations, predicted return, variance, volatility, concentration, minimum and
maximum weights, and finite-value flag.
\textbf{Checks:} compatible dimensions and variance not materially negative.

\paragraph{\texttt{directional\_optimality\_check(weights, objective, lower,
upper, n\_directions, step, seed)}}
\textbf{Inputs:} feasible weights, callable objective, bounds, and diagnostic
settings. \textbf{Work:} generate reproducible budget-preserving directions,
retain feasible perturbations, and evaluate objective changes.
\textbf{Return:} number tested, smallest change, fraction improving, and seed.
\textbf{Interpretation:} a detected improvement rejects numerical optimality;
no detected improvement is not a proof.

\subsection*{Stage 4: model selection}

\paragraph{\texttt{select\_covariance\_hyperparameters(train, candidate\_k,
candidate\_tau, validation\_length)}}
\textbf{Inputs:} training data only and candidate grids.
\textbf{Work:} make a chronological estimation/validation split; refit every
candidate on the estimation block; build GMV weights; evaluate realized
validation volatility without re-estimating on validation.
\textbf{Return:} one row per candidate with estimator name, hyperparameter,
date boundaries, solver status, realized risk, concentration, and selection
rank. \textbf{Checks:} no test-data argument, nonempty blocks, and deterministic
tie-breaking.

\subsection*{Stage 5: backtesting}

\paragraph{\texttt{drift\_weights(post\_trade\_weights,
realized\_asset\_return)}}
\textbf{Inputs:} two $(n,)$ arrays. \textbf{Work:} apply the self-financing
drift formula after one realized return. \textbf{Return:} next pre-trade
weights. \textbf{Checks:} positive portfolio gross return, finite output,
budget preservation, and unchanged weights when all asset returns are equal.

\paragraph{\texttt{performance\_metrics(net\_returns, turnover, costs,
risk\_free\_rate, annualization)}}
\textbf{Inputs:} aligned dated series and annual risk-free rate.
\textbf{Work:} construct wealth and drawdown; compute compounded annual return,
annual volatility, Sharpe, maximum drawdown, average rebalance turnover, total
cost, and final wealth. \textbf{Return:} a scalar dictionary.
\textbf{Reject:} empty, misaligned, duplicated-date, or nonfinite inputs.

\paragraph{\texttt{frozen\_weight\_backtest(test, weights,
charge\_initial\_trade, cfg)}}
\textbf{Inputs:} test returns and one weight vector chosen from train.
\textbf{Work:} apply constant target weights to every test date, follow the
declared initial-cost convention, and compute all paths.
\textbf{Return:} dictionary with gross/net returns, costs, wealth, drawdown,
weights, and metrics. \textbf{Check:} weights never depend on test observations.

\paragraph{\texttt{backtest\_rebalanced\_strategy(train, test,
covariance\_estimator, optimizer, cfg, window)}}
\textbf{Inputs:} separated samples, two callables, configuration, and
\texttt{rolling} or \texttt{expanding}. \textbf{Work:} at each date build an
information set ending at $t-1$, rebalance on schedule, compute turnover from
pre-trade weights, apply post-trade weights to $r_t$, deduct costs, then drift.
\textbf{Return:} aligned target and pre-trade weights, gross/net returns,
turnover, costs, wealth, drawdown, solver logs, estimation-window dates, and
summary metrics. \textbf{Checks:} one-period lag, feasible targets, budget after
drift, and explicit handling of solver failure.

\paragraph{\texttt{plot\_backtest\_diagnostics(backtests, output\_dir)}}
\textbf{Input:} dictionary of strategy backtest outputs.
\textbf{Work:} align dates and create wealth, drawdown, rolling-volatility,
turnover, and weight-heat-map figures with consistent colors.
\textbf{Return:} \texttt{None}. \textbf{Checks:} labeled axes and units,
legends, nonempty files, and closed figures.

\subsection*{Stage 6: orchestration and tests}

\paragraph{\texttt{run\_one\_seed(seed)}}
\textbf{Input:} one integer seed. \textbf{Work:} run the same generation,
selection, and backtest pipeline, reselecting hyperparameters inside that
seed's training sample. \textbf{Return:} one tidy data frame with seed,
strategy, selected parameters, and all metrics. \textbf{Rule:} never copy a
choice made using another seed's test result.

\paragraph{\texttt{run\_sanity\_checks()}}
\textbf{Input:} none. \textbf{Work:} implement the deterministic checks listed
in the skeleton, including identity-covariance GMV, PCA reconstruction, weight
drift, invalid covariance rejection, and one-period lag.
\textbf{Return:} \texttt{None}; raise immediately on failure and print one
success message only after every check passes.

\paragraph{\texttt{main()}}
\textbf{Input:} none; read settings from \texttt{CFG}.
\textbf{Work:} create directories, run sanity checks, generate or load data,
fit diagnostics, select hyperparameters, construct portfolios, run backtests,
save tables and plots, and execute multi-seed robustness.
\textbf{Return:} \texttt{None}. \textbf{Requirement:} one top-to-bottom run
from an empty data/output directory should reproduce every generated artifact.

\end{sloppypar}

\section{Stage 0 -- Set up a reproducible experiment}

Use \texttt{seed=2026}, $n=10$ assets, $k=3$ factors,
$T_{\mathrm{train}}=756$, and $T_{\mathrm{test}}=504$ trading days.

\begin{enumerate}[label=\textbf{0.\arabic*.}]
\item Create a configuration object containing all seeds, dimensions,
annualization constants, transaction costs, and file paths. Avoid unexplained
numbers inside functions.
\item Create separate random-number generators for population parameters,
training returns, and test returns. Explain why separate streams make debugging
easier.
\item Write a validation function that rejects non-finite arrays, unexpected
dimensions, asymmetric covariance matrices, and materially negative
eigenvalues.
\item Run the empty skeleton once and confirm that missing blocks fail with an
explicit \texttt{NotImplementedError}, not a cryptic downstream error.
\end{enumerate}

\begin{checkpointbox}{Self-check 0}
Your report states the software versions, random seed, annualization convention,
and directory structure. Running the project twice should produce identical data.
\end{checkpointbox}

\section{Stage 1 -- Generate train and test returns}

Generate returns from the factor model
\[
r_t=\alpha+Bf_t+\varepsilon_t,
\]
where
\[
f_t\sim\mathcal N\!\left(\frac{\mu_f}{252},
\frac{\Sigma_f}{252}\right),\qquad
\varepsilon_t\sim\mathcal N\!\left(0,\frac{D}{252}\right).
\]
Use annual factor means
$\mu_f=(0.05,0.025,0.015)\trans$, annual factor volatilities
$(0.16,0.10,0.07)$, and
\[
R_f=\begin{pmatrix}
1&0.25&-0.10\\0.25&1&0.20\\-0.10&0.20&1
\end{pmatrix},
\qquad
\Sigma_f=\diag(0.16,0.10,0.07)R_f\diag(0.16,0.10,0.07).
\]
Draw $B_{ij}\sim\mathcal N(0,0.45^2)$, annual alphas uniformly between
$-1\%$ and $3\%$, and annual idiosyncratic volatilities uniformly between
$10\%$ and $22\%$. If these volatilities are $s_i$, set
$D=\diag(s_1^2,\ldots,s_n^2)$. Use the same population parameters for train
and test, but independent shocks.

\begin{enumerate}[label=\textbf{1.\arabic*.}]
\item Construct $\Sigma_f$ from volatilities and correlations. Check symmetry,
diagonal entries, and smallest eigenvalue before simulating anything.
\item Generate $\alpha$, $B$, and $D$. Record their shapes and units. Which
quantities are annual and which are daily?
\item Derive and compute the theoretical annual moments
\[
\mu_{\mathrm{true}}=\alpha+B\mu_f,
\qquad
\Sigma_{\mathrm{true}}=B\Sigma_fB\trans+D.
\]
\item Simulate train and test shocks independently. Verify that the same
population parameters, not newly drawn parameters, are used for both samples.
\item Attach non-overlapping business-day indices. Name columns
\texttt{Asset\_01}, \ldots, \texttt{Asset\_10} and save exactly two return
CSV files.
\item Save the population parameters in \texttt{true\_parameters.npz}. Reload
all three files and verify exact shapes, ordering, and absence of missing data.
\item Compare annualized sample means and covariances with theoretical values.
Report relative Frobenius covariance error
$\norm{\widehat\Sigma-\Sigma_{\mathrm{true}}}_F/
\norm{\Sigma_{\mathrm{true}}}_F$ for train and test.
\item Explain why the test estimates need not be closer to the truth even
though the data-generating process is unchanged.
\end{enumerate}

\begin{checkpointbox}{Self-check 1}
Provide a table with sample sizes, missing-value counts, minimum covariance
eigenvalue, mean-estimation error, and covariance-estimation error for each
sample. Do not continue if units or dates are inconsistent.
\end{checkpointbox}

\section{Stage 2 -- Learn covariance structure with PCA}

Fit PCA only on centered training returns. You may use an eigendecomposition
of the training covariance or an SVD of the centered data matrix. Never fit a
transformation on the test sample.

\begin{enumerate}[label=\textbf{2.\arabic*.}]
\item Compute the training mean, center the training matrix, and verify that
each centered column has mean numerically close to zero.
\item Compute eigenvalues and eigenvectors in decreasing order. Verify
$U\trans U\approx I$ and reconstruct the empirical covariance from all
components.
\item Plot eigenvalues and cumulative explained variance. Mark $k=3$. What
feature of the data-generating process motivates three components?
\item Construct the rank-three and diagonal-corrected PCA covariances. Check
symmetry, diagonal preservation, rank, and smallest eigenvalue.
\item Compare empirical and PCA covariances with
$\Sigma_{\mathrm{true}}/252$ using Frobenius error. Separate approximation
error from estimation error in your discussion.
\item Compare empirical and theoretical leading subspaces. If
$s_j$ are singular values of $U_{\mathrm{true}}\trans U_{\mathrm{PCA}}$,
use principal angles $\theta_j=\arccos(\min(1,\max(-1,s_j)))$. Explain why
eigenvector signs should not be compared directly.
\item Transform train and test returns using the training mean and loadings.
Compare score variances without refitting PCA on test data.
\item Explain why PCA components need not recover the columns of $B$ exactly,
especially when idiosyncratic variances differ across assets.
\end{enumerate}

\begin{checkpointbox}{Self-check 2}
Suggested figures: scree plot, cumulative explained variance, empirical versus
theoretical eigenvalues, and a heat map of the diagonal-corrected covariance.
Suggested table: explained variance, principal angles, covariance errors, and
condition numbers.
\end{checkpointbox}

\section{Stage 3 -- Construct and validate portfolios}

From training returns compute annualized $\widehat\mu$ and
$\widehat\Sigma$. Build the following estimators: empirical covariance,
diagonal-corrected PCA covariance for several $k$, and ridge covariance for a
grid of $\tau$ values.

\begin{enumerate}[label=\textbf{3.\arabic*.}]
\item Implement a portfolio validator reporting budget error, minimum weight,
maximum weight, predicted return, predicted volatility, concentration, and
all stated constraint violations.
\item Implement equal weight and the closed-form unconstrained GMV portfolio.
Check the closed-form stationarity equation before using an optimizer.
\item Solve long-only GMV with $0\leq w_i\leq0.30$. Report solver status and
identify active lower and upper bounds using a numerical tolerance.
\item Solve long-only maximum Sharpe with $r_f=2\%$. Guard against a zero-risk
denominator and run several feasible initializations to detect local numerical
issues.
\item Solve bounded mean--variance problems
\[
\min_w\ \frac12w\trans\widehat\Sigma w-\gamma\widehat\mu\trans w,
\quad \mathbf1\trans w=1,\quad0\leq w_i\leq0.30
\]
for at least six logarithmically spaced values of $\gamma$.
\item Trace a target-return efficient frontier. Before solving, determine a
feasible target interval under the weight bounds; do not ask the solver to
reach an impossible return.
\item Compare empirical, PCA, and ridge covariances through eigenvalues,
condition numbers, weights, concentration, predicted risk, and solver
residuals.
\item Write down the KKT conditions for long-only GMV. If the solver exposes
multipliers, compute stationarity and complementary-slackness residuals;
otherwise verify primal feasibility and a directional optimality check.
\end{enumerate}

\begin{checkpointbox}{Self-check 3}
No portfolio enters the backtest unless its weights are finite, budget error is
below $10^{-8}$, bound violations are below $10^{-8}$, predicted variance is
nonnegative up to tolerance, and the solver reports success.
\end{checkpointbox}

\section{Stage 4 -- Choose hyperparameters without test leakage}

The number of PCA components $k$, ridge parameter $\tau$, and mean--variance
parameter $\gamma$ are hyperparameters. Choose them using training information
only.

\begin{enumerate}[label=\textbf{4.\arabic*.}]
\item Reserve the final 126 training observations as a chronological validation
block. Estimate each candidate using only the earlier training observations.
\item State a selection criterion before examining validation results. A
reasonable primary criterion is realized validation volatility, with average
turnover or concentration as a tie-breaker.
\item Compare at least $k\in\{1,2,3,5\}$ and
$\tau\in\{10^{-4},10^{-3},10^{-2},10^{-1},1\}$. Explain why selecting the
smallest in-sample predicted variance would favor overfitting.
\item Freeze the chosen rule before loading test returns. Store the selected
parameters and criterion in a results table.
\item Optional: replace the single validation block with expanding or rolling
time-series validation and compare the chosen hyperparameters.
\end{enumerate}

\section{Stage 5 -- Backtest without look-ahead bias}

Begin with a frozen-weight experiment: estimate once on the entire training
sample, then hold the weights throughout the test sample. Next implement a
rebalanced strategy.

At decision date $t$, use only observations strictly before $t$. Let $w_t^-$
be pre-trade weights and $w_t^+$ post-trade weights. Between trades, weights
drift according to
\[
w_{i,t}^-=
\frac{w_{i,t-1}^+(1+r_{i,t-1})}
{1+(w_{t-1}^+)\trans r_{t-1}}.
\]
At a rebalance,
\[
\operatorname{TO}_t=\sum_i|w_{i,t}^+-w_{i,t}^-|,
\qquad
r_{p,t}^{\mathrm{net}}=(w_t^+)\trans r_t-c\operatorname{TO}_t,
\]
with $c=0.001$. At a non-rebalance date set $w_t^+=w_t^-$ and
$\operatorname{TO}_t=0$.

\begin{enumerate}[label=\textbf{5.\arabic*.}]
\item Write a one-period-lag test: changing return $r_t$ must not change the
weights applied to $r_t$. Include this test in the code.
\item Rebalance every 21 trading days. Initially use the final 252 training
observations; later use the most recent 252 observations available before the
decision date.
\item Compare frozen weights, a 252-day rolling window, and an expanding
window. State precisely which dates enter each estimate.
\item Before every trade, update drifting weights and verify they still sum to
one. Then compute turnover relative to the new target weights.
\item Choose and report the initial-trade convention. Either begin from equal
weight and charge turnover, or treat initial allocation as costless; use the
same convention for every strategy.
\item Compute wealth $V_t=V_{t-1}(1+r_{p,t}^{\mathrm{net}})$ with $V_0=1$ and
drawdown $1-V_t/\max_{s\leq t}V_s$.
\item Report annualized return, volatility, Sharpe ratio, maximum drawdown,
average turnover on trade dates, total transaction costs, and final wealth.
Use daily risk-free return $r_{f,d}=(1+r_f)^{1/252}-1$.
\item Compare equal weight and empirical-, PCA-, and ridge-covariance GMV.
Attribute differences separately to estimated risk, realized risk,
concentration, turnover, and costs.
\end{enumerate}

\begin{checkpointbox}{Self-check 5}
Suggested diagnostics: wealth, drawdown, rolling 63-day volatility, turnover,
weight heat maps, gross/net performance, and evidence that the first test
weight excludes the first test return.
\end{checkpointbox}

\section{Stage 6 -- Robustness and interpretation}
\enlargethispage{\baselineskip}

\begin{enumerate}[label=\textbf{6.\arabic*.}]
\item Use the non-deployable oracle GMV portfolio from
$\Sigma_{\mathrm{true}}$ only to quantify estimation error.
\item Repeat the experiment for at least five additional seeds, reselecting
hyperparameters within each training sample. Plot the distributions of realized
volatility, Sharpe ratio, drawdown, turnover, and final wealth; report medians
and dispersion rather than means alone.
\item Across seeds, decide whether PCA or ridge consistently improves
conditioning, weight stability, or realized risk. Discuss limitations from
Gaussian returns, constant parameters, mean-estimation error, and simplified
transaction costs.
\end{enumerate}

\end{document}
